\section{Montague's generalized semantics: reconstruction proofs}
\label{sec:annex-p08-n-b-006-02}

The following proofs make explicit the semantic arguments stated or
indicated in the Synthese presentation. They are reconstructions from
its definitions, not claims that it printed these full derivations.
The compactness, effectiveness, and natural-axiomatization assertions
retain their cited boundaries in
\cref{thm:p08-montague-metatheory-source}
\cite{Montague1970Pragmatics}.

\begin{annexitem}{annex:p8.n-b-006.2.pragmatic}{thm:p08-montague-pragmatic-recursion}{The simultaneous recursion and assignment lemma}
\begin{proof}
Order pairs \((i,\phi)\) by proper subformula in the second coordinate,
allowing arbitrary change of index in the first. There is no infinite
descending chain, because formula length decreases at every step. In an
operator clause, every required \((j,\psi)\) precedes \((i,N\psi)\).
Thus all formula extensions are uniquely determined by well-founded
recursion. Equivalently one may construct the entire function of the
index at each formula-complexity stage; this is ordinary structural
recursion on intensions.

Suppose two assignments agree on all free variables of \(\phi\).
Atoms agree by their displayed coordinate evaluations. Boolean
operations preserve agreement. For a quantified variable, matching
replacements by the same \(y\in U\) agree on the matrix's free
variables. For an operator, the induction hypothesis applies at every
index to the same two assignments, so their truth sets are identical
and the same relation \(R_N\) decides truth. Induction proves the
assignment lemma. With no free variables, either all assignments or
none are in the extension.

For \(U\ne\varnothing\), constant sequences and coordinate replacements
exist. Evaluating the first four criteria is therefore exactly evaluating
the atomic clauses; existence of an assignment with the indicated
coordinate is equivalent to the respective displayed atomic condition.
Complement and intersection give the next two criteria. The quantifier
clause gives the seventh: a witness coordinate supplies satisfaction,
and a satisfying assignment supplies that witness. The sentence
assignment lemma makes each set
\(\{j:x\in\operatorname{Ext}_j(\phi)\}\) independent of \(x\),
and equal to its set of true indices. This gives the eighth criterion.
If \(U\) is empty, extensions equal the empty assignment universe and
sentences are all true; such interpretations cannot falsify a conclusion.
Removing them therefore leaves sentence consequence unchanged.
\end{proof}
\end{annexitem}

\begin{annexitem}{annex:p8.n-b-006.2.examples}{thm:p08-montague-operator-examples}{Reading the temporal and relational clauses}
\begin{proof}
For a sentence \(\phi\) let \(J\) be its truth set. With the past
relation the clause says \(\exists j\in J\,(j<i)\), which is
precisely truth at some earlier time; replacing the strict order by its
reverse gives the future case. With \(R_N(i,J)\) defined by
\(r[i]\subseteq J\), it says truth at every accessible point.
For the progressive relation it says exactly that some interval
\((a,b)\) with \(a<i<b\) is included in \(J\). This is equivalent
to truth throughout that interval. The set-of-indices relation can
express this without a designated set of individual successor indices;
the proposition set itself is its semantic argument. In fact no fixed
successor set \(T_i\) represents this family by the universal relational
clause. Every open interval about \(i\) would have to contain
\(T_i\), so \(T_i\subseteq\{i\}\). If it is empty, that relational
clause accepts the empty truth set; if it is \(\{i\}\), it accepts
the singleton truth set. Neither contains an interval about \(i\).
Nor can the existential relational clause represent it: any nonempty
successor set accepts some singleton, and an empty successor set rejects
the entire real line. This verifies the generalized-semantics motivation
under the standard two relational clauses. Probability-threshold
operators and conditional necessity are other motivations mentioned by
the source, but no particular probability measure or conditional
construction is fixed by that mention.
\end{proof}
\end{annexitem}

\begin{annexitem}{annex:p8.n-b-006.2.intensional}{thm:p08-montague-intensional-recursion}{Typed recursion, descriptions, and the five new criteria}
\begin{proof}
Use expression complexity, simultaneously at all indices, for terms and
formulas. A description refers only to its proper matrix formula, albeit
at each altered predicate coordinate. Applications refer to proper
argument terms; Boolean, quantifier, and \(\Box\) clauses refer to
proper formulas. Hence the simultaneous recursion is well founded.
The full domains \(\mathcal P_k\) are sets. The matrix-satisfying
candidate set of a description is a subset of \(\mathcal P_k\).
Choosing its unique member, or the fixed member \(O_k\), is therefore
a well-defined, total term evaluation of the correct sort.

The free-variable induction extends the preceding proof. The extra
case for a description compares its candidate sets after changing its
bound coordinate to the same arbitrary \(Y\); they are equal by
induction, so either both choose the same unique member or both choose
the same default. The \(\Box\) case uses agreement at every index.
Thus sentences again have full or empty extension.

An existential predicate witness is exactly a coordinate value from
\(\mathcal P_k\), proving criterion 8'. The typed constant-application
clause tests \(\langle Y\rangle\in F_P(i)\) with the whole assigned \(Y\), not
\(Y(i)\), proving 9'. The universal-index clause proves 10'. The
description clause and the typed application clause together give the
two disjoint alternatives in 11': one unique matrix-satisfying predicate
accepted by \(F_P(i)\), or failure of uniqueness with \(O_k\) accepted
there. Finally the zero-place application tests membership of the empty
tuple in \(Y(i)\). The only two zero-place relations are \(\mathrm F\)
and \(\mathrm T\), proving 12'. The seven nonmodal criteria follow
as before, with the typed assignment set in place of \(U^\omega\).
\end{proof}
\end{annexitem}

\begin{annexitem}{annex:p8.n-b-006.2.concepts}{thm:p08-montague-concept-relations}{Singleton intensions and constant relations}
\begin{proof}
The first formula says at every index that there is exactly one
\(a\in U\) satisfying the unary predicate. It therefore says
\(Y(i)=\{f(i)\}\) for a uniquely determined function \(f:I\to U\).
Conversely any such function gives the predicate \(i\mapsto\{f(i)\}\)
and satisfies the formula. This proves the claimed bijection, rather
than only an injection into predicates.

For any tuple \((a,b)\in U^2\), the second formula says either it
belongs to \(Y(j)\) at every index \(j\), or it belongs at no index.
Hence all \(Y(j)\) equal one relation \(T\subseteq U^2\).
Conversely a constant function \(j\mapsto T\) gives the indicated
disjunction for each tuple. The argument with \(k\) universally
quantified individual variables proves the same statement at arity
\(k\). At arity zero it is \(\Box G[\ ]\lor\Box\neg G[\ ]\),
which singles out the constantly true and constantly false propositions.
\end{proof}
\end{annexitem}

\begin{annexitem}{annex:p8.n-b-006.2.descriptions}{thm:p08-montague-descriptions}{Eliminating descriptions and proving the name construction}
\begin{proof}
For same-arity predicate variables define, as an abbreviation,
\[
 \operatorname{Eq}_k(G,H)=\Box\forall u_0\cdots u_{k-1}
                   (G[\bar u]\leftrightarrow H[\bar u]),\qquad
 \operatorname{Empty}_k(G)=\Box\forall\bar u\,\neg G[\bar u].
\]
For \(k=0\) omit the individual quantifiers and arguments. Full
predicate functions are equal iff their values agree on all tuples at
all indices, so these abbreviations test actual equality of intensions
and the unique empty intension, respectively.

Rename bound variables to avoid capture. For a description
\(\mathsf T H\psi\), with \(H\) of arity \(k\), write
\(\psi(G)\) for substitution of a fresh variable \(G\) for its
free \(H\)-occurrences. Put
\begin{align*}
 \operatorname{Unique}_\psi(G)&=
   \psi(G)\land\forall K(\psi(K)\to\operatorname{Eq}_k(K,G)),\\
 \operatorname{Select}_\psi(G)&=
   \operatorname{Unique}_\psi(G)\lor
   (\neg\exists K\operatorname{Unique}_\psi(K)
                        \land\operatorname{Empty}_k(G)).
\end{align*}
At any index and assignment exactly one \(G\)-value satisfies
\(\operatorname{Select}_\psi\): the unique candidate if there is
one, or \(O_k\) if there is not. Replace an atomic formula
\(\Theta[\mathsf T H\psi]\), including the case where the term
occurs in predicate position, by
\(\exists G(\operatorname{Select}_\psi(G)\land\Theta[G])\).
It has the same truth, since it quantifies precisely the value of the
description. For several occurrences use distinct fresh variables.
Eliminate descriptions in the matrices first, then in their containing
atoms, and extend compositionally through connectives, quantifiers,
and \(\Box\). Every replacement is pointwise equivalent at every
index and assignment, so equivalence remains valid inside a modality.
The number and nesting depth of descriptions decrease, yielding a
description-free sentence after finitely many steps.

In the source's displayed unary example \(P[\mathsf T G Q[G]]\),
the result can be written
\[
 \exists G(Q[G]\land\forall H(Q[H]\to\operatorname{Eq}_1(H,G))
                    \land P[G])\lor
 (\neg\exists G(Q[G]\land\forall H(Q[H]\to\operatorname{Eq}_1(H,G)))
             \land\exists G(\operatorname{Empty}_1(G)\land P[G])).
\]
This keeps both the unique-candidate and default alternatives on source
p.~82, rather than silently treating descriptions as partial terms.

For a name \(\widehat{\bar u}\phi\), fix the remaining coordinates
of an assignment and define \(Y(j)\) by the displayed comprehension
in the body. It is a member of the full \(\mathcal P_m\). At all
tuples and all indices it agrees with \(\phi\); therefore it satisfies
the description's defining formula. Any other candidate must agree at
every such tuple and index, so is the same function. The selected
predicate is thus always that unique \(Y\), never the default from
failed uniqueness. At arity zero its truth value at \(j\) is exactly
the truth of \(\phi\) with the fixed remaining assignment. Equality
of this proposition with the everywhere-true proposition is equivalent
to the universal-index truth condition for \(\Box\phi\).

For the iterated-belief example, fix \(x\). The inner proposition
\(\widehat{S[x]}\) is the unique \(H\) satisfying
\(\Box(H[\ ]\leftrightarrow S[x])\). Consequently
\(B[R,\widehat{S[x]}]\) is pointwise equivalent to
\(\exists H(B[R,H]\land\Box(H[\ ]\leftrightarrow S[x]))\).
The outer name is then the unique \(G\) whose universal-index truth
agrees with this formula. Applying \(B[J,G]\), and preserving the
outer actuality restriction and individual existential, gives exactly
the displayed description-free belief sentence. If sentences \(\phi\)
and \(\psi\) are logically equivalent, their truth functions agree
at every index of each interpretation. Their names therefore denote
the same proposition, and every empirical belief relation accepts
one precisely when it accepts the other. This proves the stated
semantic equivalence without introducing psychological axioms.
\end{proof}
\end{annexitem}

\begin{annexitem}{annex:p8.n-b-006.2.reduction}{thm:p08-montague-pragmatic-reduction}{Truth preservation in the forward translation}
\begin{proof}
For nonempty \(U\), every ordinary individual assignment extends to a
typed system: assign arbitrary predicate functions to its other
coordinates. Such functions always exist, for example the empty ones.
The translation has no free predicate variables, since every introduced
\(G\) is freshly bound. Inductively its truth depends only on the
original individual coordinates. Original atoms agree because their
interpretations are unchanged; Boolean clauses agree; individual
quantifiers range over the same \(U\).

Fix the individual assignment \(x\), and let
\(J=\{j:x\in\operatorname{Ext}_{j}(\psi)\}\). By induction this
is also the translated formula's truth set. There is exactly one
\(Y\in\mathcal P_0\) with \(Y(j)=\mathrm T\) iff \(j\in J\).
The formula \(\Box(G[\ ]\leftrightarrow\psi^*)\) restricts
\(G\) to this \(Y\), since \(G\) is absent from \(\psi^*\).
The other conjunct then holds iff \(\langle Y\rangle\in F'_{N'}(i)\), which by
definition is iff \((i,J)\in R_N\). This is precisely the pragmatic
operator clause. Structural induction proves truth preservation for
formulas with fixed assignments and therefore for sentences. If
\(U\) is empty, both assignment spaces are empty and both sentence
truth definitions make every sentence true, so the sentence claim
continues to hold in that boundary case.
\end{proof}
\end{annexitem}

\begin{annexitem}{annex:p8.n-b-006.2.converse}{thm:p08-montague-converse-reduction}{Truth-set bijection and the partial reverse translation}
\begin{proof}
Send a proposition \(Y\) to \(J_Y=\{j:Y(j)=\mathrm T\}\).
Its inverse sends \(J\subseteq I\) to the function with value
\(\mathrm T\) on \(J\) and \(\mathrm F\) off \(J\).
Thus this is a bijection, including the empty index-set case as a
set-theoretic correspondence. At each index, subsets of the proposition
domain correspond bijectively to subsets of \(\mathcal P(I)\).
Consequently the two displayed predicate/modality constructions are
inverse, not merely one-sided encodings.

For a source subformula \(P[\widehat\psi]\), the name theorem gives
the proposition with truth set \(J_\psi\). Its typed application
holds iff that proposition belongs to the interpretation of \(P\),
which by the constructed \(R_{P'}\) is iff
\((i,J_\psi)\in R_{P'}\). This is the target's \(P'\psi^*\).
The primitive \(\Box\) likewise has its pragmatic representation
because truth at every index means exactly \(J_\psi=I\).
Atoms not involved in replacements keep their interpretations, and
connectives and individual quantifiers commute with the translation.
Induction therefore proves preservation for exactly the sublanguage
whose translated formula is pragmatic. A free or quantified predicate
variable with no corresponding replacement is outside this induction's
grammar; the conclusion is not asserted for it. Empty interpretations
are harmless for the sentence statement as in the forward direction.
\end{proof}
\end{annexitem}

\begin{annexitem}{annex:p8.n-b-006.2.extended}{thm:p08-montague-extended-reduction}{Binding the arguments' predicate intensions}
\begin{proof}
With \(m=p=0\) the operator clause tests only an \(n\)-tuple in
\(F_N(i)\), exactly the predicate-constant clause. With \(m=n=0,p=1\)
its only argument is the proposition whose truth set is the operand's
truth set. Given an ordinary \(R_N\), interpret that argument as
accepted iff the corresponding pair \((i,J)\) is in \(R_N\).
The inverse follows from the proposition/truth-set bijection.

For arbitrary degree, use the same \(A\), the same individual
constant interpretations, and let the new typed predicate's relation
at \(i\) equal the original \(F_N(i)\). At assignment \(x\), the
name \(\widehat{\bar u}\phi_q^*\) denotes exactly the function
\(Y_{q,x}\), by the name theorem and induction on \(\phi_q\).
Its typed application therefore tests exactly the same tuple of
individuals and predicate functions as the extended operator clause.
The other extension clauses agree compositionally. This proves truth
preservation, including proper binding of the displayed \(u_r\).

Belief supplies one individual and one proposition, giving degree
\(\langle0,1,1\rangle\). Seeming supplies two individuals and one
unary-property function constructed by binding one individual variable,
giving degree \(\langle1,2,1\rangle\). To preserve argument order,
interpret the target typed relation on \((u,v,Y)\) by the original
seeming relation on \((u,Y,v)\). These checks establish the displayed
examples. They introduce predicate names only for specific formula
arguments, not variables ranging over arbitrary predicates. Thus the
translation gives the claimed inclusion, not a full reverse reduction
of predicate quantification.
\end{proof}
\end{annexitem}

\begin{annexitem}{annex:p8.n-b-006.2.substitution}{thm:p08-montague-substitution}{Rigid variable assignments and a two-index counterexample}
\begin{proof}
An assigned individual is fixed while the index changes. If \(u,v\)
have equal assigned values, their atomic occurrences evaluate equally
at every index. The usual capture-free substitution induction passes
through connectives and quantifiers. It also passes through ordinary
operators, because their argument truth sets agree at every index,
through \(\Box\), because each indexed truth agrees, and through
descriptions, because the candidate sets agree. In extended pragmatics
the same induction compares each bound-variable predicate function
\(Y_{q,x}\). Hence replacement of a free variable occurrence preserves
truth in all three systems. Quantifying the two variables proves the
stated validity. Instantiation by a variable is the same assignment
replacement argument.

For constants let \(I=\{0,1\}\), \(U=\{a,b\}\), and both
\(A_i=U\). Let \(c(0)=d(0)=a\), but \(c(1)=a,d(1)=b\), and
interpret an ordinary predicate \(P\) as \(\{a\}\) at both indices.
At index zero \(c=d\) is true. Nevertheless \(\Box P[c]\) is true
and \(\Box P[d]\) is false. In ordinary pragmatics represent the
same universal-index operator by \(R_N(i,J)\) iff \(J=I\); in
extended pragmatics use its degree-\(\langle0,0,1\rangle\) encoding.
Thus the failure occurs in all three systems.

For failed constant instantiation, take \(I=\{0,1\}\),
\(U=\{a,b\}\), both actuality sets full, \(P(0)=\{a\}\),
\(P(1)=\{b\}\), and \(c(0)=a,c(1)=b\).
At index zero every rigid individual fails \(\Box P[v]\):
\(a\) fails at one and \(b\) fails at zero. Hence
\(\forall v\neg\Box P[v]\) is true there. But \(\Box P[c]\)
is true, so its proposed constant instance \(\neg\Box P[c]\) is
false. The same encodings give the corresponding pragmatic examples.
\end{proof}
\end{annexitem}

\begin{annexitem}{annex:p8.n-b-006.2.neighborhood}{thm:p08-montague-neighborhood-bridge}{The exact relation/function correspondence}
\begin{proof}
For every \(i\) and \(J\subseteq I\), the definition gives
\(J\in\mathcal N_N(i)\) iff \((i,J)\in R_N\). Substituting the
sentence's truth set into this equivalence yields exactly the modern
neighborhood clause. Conversely, from a neighborhood function form
\(R_N=\{(i,J):J\in\mathcal N_N(i)\}\). Taking sections returns
the original function, and taking their union returns the original
relation.

For propositional valuations use zero-place predicate constants: a
constant is true at \(i\) iff its empty tuple lies in its extension.
Arbitrary sets of true indices give precisely arbitrary propositional
valuations. Atoms therefore agree, Boolean clauses agree, and the
displayed equivalence makes the operator step agree. Induction gives
truth preservation for every propositional modal formula. Choosing
\(U\) to be a singleton and every \(A_i=U\) realizes any such
neighborhood model inside Montague's nonempty interpretations. No
additional closure is introduced by this construction.
\end{proof}
\end{annexitem}

\begin{annexitem}{annex:p8.n-b-006.2.modern}{thm:p08-montague-modern-boundary}{Modern closure names and their separating examples}
\begin{proof}
For a fixed valuation write \(X,Y\) for proposition truth sets.
Upward closure gives \(X\cap Y\in\mathcal N(i)\Rightarrow
X,Y\in\mathcal N(i)\), which is M. Binary-intersection closure
gives \(X,Y\in\mathcal N(i)\Rightarrow X\cap Y\in\mathcal N(i)\),
which is C. Containing \(I\) gives \(\Box\top\), which is N.
Conversely M forces upward closure under arbitrary valuations: for
\(X\subseteq Y\), value its two atoms by \(X,Y\), so their
intersection is \(X\). The same arbitrary-set valuations make C
force intersection closure. Equivalent formulas have equal truth sets
in every model, so replacement follows without any of these closures.

On a nonempty set let every neighborhood family be empty. It is upward
and intersection closed, but not unit containing, so EMC alone does
not imply N. For EN without monotonicity, take
\(I=\{0,1,2\}\) and at index zero the family
\(\{I,\{0\}\}\): it contains the unit and \(\{0\}\), but not
the superset \(\{0,1\}\). For EN without intersection closure,
take \(I=\{0,1\}\) and the family
\(\{I,\{0\},\{1\}\}\). The intersection of the two singleton
members is absent. Indeed let \(p\) have truth set \(\{0\}\) and
\(q\) have empty truth set. At zero both \(\Box(p\to q)\) and
\(\Box p\) are true but \(\Box q\) is false. This directly
refutes K on a unit-containing neighborhood model.

At theoremhood the classical congruential base plus M derives modal
monotonicity for any valid implication \(a\to b\): classical
equivalence of \(a\) with \(a\land b\), replacement, and M give
\(\Box a\to\Box b\). Applying C to
\(\Box(a\to b)\land\Box a\) gives
\(\Box((a\to b)\land a)\); monotonicity gives \(\Box b\).
This derives K. N and replacement give theorem-only necessitation:
if \(a\) is a theorem it is equivalent to \(\top\), whose box is
available. Conversely normal K with theorem-only necessitation derives
M from the tautological projections of \(a\land b\), C from the
tautology \(a\to(b\to a\land b)\), and N by necessitating
\(\top\). Necessitating the two directions of a theorem-level
equivalence and applying K gives replacement. This proves the claimed
theoremhood equivalence of the named packages, without asserting a
new strong-completeness result.

Finally a premise \(p\) can be true at one of two indices and false
at the other. The source's universal-index \(\Box p\) is then false
at the first index, so \(p\) does not locally entail \(\Box p\).
Necessitation of theorems must not be extended to arbitrary locally
true premises. The arbitrary pragmatic \(N\)-relations impose no
condition preventing this distinction.
\end{proof}
\end{annexitem}
